\section*{РОЗДІЛ 3. ТЕХНІЧНІ МОЖЛИВОСТІ GENAI ДЛЯ АВТОМАТИЗАЦІЇ НАУКОВОГО АДМІНІСТРУВАННЯ}

\subsection*{3.1. Обробка бібліометричних даних великих обсягів}

Цифровізація наукової комунікації призвела до стрімкого зростання обсягів бібліометричних даних і ускладнення їх структури (метадані публікацій, профілі авторів, афіліації, посилання, показники доступу тощо). У цих умовах великі мовні моделі (LLM) та суміжні інструменти обробки наукових текстів формують технологічну основу для переходу від трудомісткої ручної звітності до більш автоматизованих, відтворюваних і масштабованих процедур наукового адміністрування у ЗВО.

\subsubsection*{3.1.1. OpenAlex як ключовий відкритий ресурс}

Одним із найбільш релевантних джерел відкритих бібліометричних метаданих є \textbf{OpenAlex} --- наступник Microsoft Academic Graph і, за наявними оцінками, найбільша відкрита платформа такого типу \cite{Priem2022OpenAlex,OpenAlex2024Docs}. Її практична цінність для управлінських задач визначається поєднанням широкого охоплення, нормалізованих сутностей (автори, інституції, джерела) та можливістю автоматизованого доступу.

Станом на 2024 рік OpenAlex акумулює:
\begin{itemize}[leftmargin=*,itemsep=0.2em]
\item понад \textbf{250 мільйонів наукових праць} (works);
\item більше \textbf{100 мільйонів авторів} (authors) з автоматизованою дисамбігуацією;
\item \textbf{250+ тисяч джерел} (sources) --- журнали, конференції та репозитарії;
\item \textbf{109 тисяч інституцій} (institutions) з нормалізованими назвами;
\item \textbf{65 тисяч концептів} (concepts) для тематичної класифікації.
\end{itemize}

Для наукового адміністрування принциповими є такі властивості платформи.

\textbf{Відкритий програмний доступ.} На відміну від комерційних індексів, OpenAlex надає безкоштовний REST API, що спрощує побудову відтворюваних процесів збору й оновлення метаданих.

\textbf{Автоматизована дисамбігуація авторів.} Використання ML-алгоритмів дозволяє пов'язувати публікації з авторськими профілями навіть у випадках, коли ORCID відсутній або заповнений неповністю.

\textbf{Зв'язування з ORCID.} ORCID може слугувати «якорем» (anchor) для зіставлення профілів і підвищення надійності ідентифікації.

\textbf{Розширені метадані.} Записи про праці включають відомості про цитування, статус відкритого доступу, афіліації авторів, концепти та мережеві зв'язки між роботами, що є критичним для аналітики та побудови показників.

Нижче наведено приклад типового запиту до OpenAlex API для отримання переліку публікацій автора за ORCID:

\begin{lstlisting}
GET https://api.openalex.org/works?filter=author.orcid:0000-0003-0242-2234&per-page=100
\end{lstlisting}

\subsubsection*{3.1.2. Технологія GROBID для витягування з PDF}

Поряд із метаданими, що надходять із зовнішніх індексів, інституції часто мають справу з первинними документами у форматі PDF (статті, препринти, тези, звіти). Для формування машинно-читаних представлень таких документів застосовують \textbf{GROBID} (GeneRation Of BIbliographic Data) --- відкриту бібліотеку, орієнтовану на автоматичне витягування структурованої бібліографічної інформації з PDF \cite{Lopez2015GROBID}.

У прикладних сценаріях наукового адміністрування GROBID, зокрема, забезпечує:
\begin{itemize}[leftmargin=*,itemsep=0.2em]
\item розпізнавання структури документа (заголовок, автори, анотація, розділи, бібліографія);
\item витягування та парсинг бібліографічних посилань;
\item ідентифікацію афіліацій авторів і їх подальше зіставлення з організаціями.
\end{itemize}

\subsubsection*{3.1.3. Інструменти для бібліометричної аналітики: роль детермінованих методів}

У практиці бібліометричної аналітики істотну частину задач (агрегація, очищення, science mapping, тематичні кластери) традиційно розв'язують детермінованими методами та відкритими інструментами. Наприклад, для science mapping та побудови тематичних структур широко використовується пакет \textbf{bibliometrix} \cite{Aria2017bibliometrix}, а для тематичного моделювання --- методи на кшталт \textbf{BERTopic} \cite{Grootendorst2022BERTopic}. У контексті ResearchIQ це важливо для розмежування ролей: LLM застосовується переважно для інтерфейсу (NL-запити, пояснення, оркестрація), тоді як ядро метрик і трансформацій має бути відтворюваним та перевірюваним.

\subsection*{3.2. Предиктивна бібліометрія та аналітика на базі GenAI}

У сучасній наукометрії GenAI розглядають як компонент, що може доповнювати традиційні підходи (feature-based моделі, time series, графові методи) шляхом (i) інтерпретації слабкоструктурованих даних, (ii) генерації пояснень для користувачів та (iii) гнучкої побудови запитів і звітів. Водночас емпірична придатність таких підходів суттєво залежить від протоколів валідації та контролю галюцинацій (див. розд.~6).

\subsubsection*{3.2.1. Предиктивна бібліометрія}

Цей напрямок фокусується на прогнозуванні патернів цитування та інших форм «видимості» публікації на основі її характеристик (журнал, авторський склад, тематика, статус відкритого доступу тощо).

\subsubsection*{3.2.2. Автоматизована оцінка досліджень}

Йдеться про генерацію оціночних резюме для наукових проєктів і публікацій, що може підтримувати експертні процедури за умови прозорих критеріїв і обов'язкової верифікації результатів.

\subsubsection*{3.2.3. Наукове картування та візуалізація}

Застосування GenAI до задач science mapping передбачає побудову карт наукових ландшафтів і виявлення тематичних кластерів, що є корисним для аналізу розвитку напрямів і міждисциплінарних зв'язків.

\subsubsection*{3.2.4. Рекомендаційні системи}

У межах рекомендаційних підходів GenAI може використовуватися для підбору журналів для публікації, ідентифікації потенційних співавторів і виявлення релевантних джерел на основі профілю дослідника та його робіт.

\subsubsection*{3.2.5. Оцінка дослідницького впливу}

Нарешті, GenAI розглядають як можливий компонент розширення традиційних метрик (H-індекс, кількість цитувань) альтернативними індикаторами впливу, зокрема через аналіз повних текстів, контекстів цитування та взаємодії з відкритими даними.

\subsection*{3.3. Мульти-агентні системи та автоматизація документообігу}

Окремий клас технологічних рішень пов'язаний із мульти-агентними підходами, де складні адміністративні процеси (збір даних, перевірка, агрегація, формування текстів і форм) розкладаються на послідовність ролей і взаємодій між спеціалізованими агентами. Це дає змогу поєднати різні інструменти (API, парсери, модулі перевірки) в єдиний керований контур.

\subsubsection*{3.3.1. PaperQA як RAG-агент для наукової літератури}

Як приклад RAG-підходу до роботи з науковими текстами можна навести \textbf{PaperQA} --- спеціалізований агент, що орієнтований на пошук і синтез відповідей на запити на основі корпусів наукової літератури \cite{ODonoghue2023PaperQA}.

\subsubsection*{3.3.2. Мульти-агентні архітектури}

Узагальнено, мульти-агентна архітектура може бути представлена набором ролей, де кожна відповідає за окремий етап життєвого циклу документа або даних (планування, збір, валідація, редагування та публікація):

\begin{longtable}{@{}p{0.18\textwidth}p{0.18\textwidth}p{0.54\textwidth}@{}}
\toprule
\textbf{Агент} & \textbf{Роль} & \textbf{Функції} \\
\midrule
\textbf{Editor} & Планування & Декомпозиція завдання, визначення стратегії \\
\textbf{Researcher} & Збір даних & Пошук у базах, API-запити, скрапінг \\
\textbf{Reviewer} & Валідація & Перевірка фактів, виявлення галюцинацій \\
\textbf{Reviser} & Доопрацювання & Корекція помилок, покращення якості \\
\textbf{Writer} & Компіляція & Синтез фінального документа \\
\textbf{Publisher} & Форматування & Оформлення за вимогами, експорт \\
\bottomrule
\end{longtable}

\subsubsection*{3.3.3. Автоматизація грантового адміністрування}

Автоматизація грантового циклу (підготовка заявки, перевірка відповідності вимогам, формування супровідних документів) є ще одним прикладом застосування агентних підходів. У прикладній площині згадуються, зокрема, такі рішення:

\textbf{Granted AI} --- система, навчена на понад 500{,}000 успішних грантових заявок.

\textbf{Grantable} --- платформа з сертифікацією SOC 2 Type 2.

\textbf{DeepRFP} --- рішення, що реалізує мульти-прохідну валідацію для зниження ризику помилок і галюцинацій.

\subsubsection*{3.3.4. Критичне обмеження: проблема галюцинацій}

Разом із розширенням можливостей GenAI зберігається фундаментальне обмеження, пов'язане з надійністю генерованих тверджень і посилань. Зокрема, у порівняльних дослідженнях якості посилань, згенерованих LLM у контексті наукового письма, повідомляються суттєві частки некоректних/вигаданих референсів (hallucinated references) \cite{Chelli2024HallucinationRates}, що робить обов'язковими процедури верифікації, контроль джерел і фіксацію provenance у практичних системах автоматизованої звітності.